\chapter{System Study}

\section{Existing System}
Current protection is distributed among browser warnings, URL blocklists, antivirus
software, email filters, and security extensions. These tools are valuable but often
operate independently, provide limited explanations, or cannot inspect all modern
browser behaviours. Blocklists can miss new campaigns, while heavyweight crawlers and
dynamic analysis are unsuitable for continuous protection on every user page.

\section{Proposed System}
BrowseBlock combines an event-driven extension with a centralized analysis service.
The extension observes supported navigation, webpage, download, notification, and
extension events. The dashboard enables deliberate scans and detailed review. The
backend coordinates model inference, static analysis, RAG/LLM services, authentication,
authorization, and persistence. This division keeps common checks lightweight while
allowing expensive analysis only when evidence or user intent justifies it.

\section{System Modules}
\begin{enumerate}
\item URL and fake-QR phishing detection
\item Cross-site scripting detection and prevention
\item Credential-harvesting detection
\item HTML/JavaScript redirect analysis
\item Clickjacking detection
\item Cryptojacking detection
\item Malicious file-download analysis
\item Email phishing detection
\item Browser notification-scam detection
\item Browser-extension abuse detection
\end{enumerate}

\section{Users and Roles}
General users need concise warnings and practical recommendations. Organizational users
need an additional protection layer for browsing, email, and downloads. Authorized
security analysts review evidence and threat reports, while administrators manage
configuration, services, and access. All role-specific dashboard data is subject to
authenticated ownership or explicit authorization.

\section{Constraints}
Detection is bounded by Chromium permissions, browser API visibility, network
availability, model quality, external service availability, and the evidence observable
at runtime. BrowseBlock therefore complements rather than replaces existing endpoint
security. Supported browsers, model versions, datasets, file limits, retention periods,
and the first-prototype module set remain implementation decisions to be finalized.
